% Môn: Vật lý
% Lớp: 12
% Chương: Chương I. Vật lí nhiệt
% Bài: L12C1 Bài 3. Nhiệt độ. Thang nhiệt độ – nhiệt kế
% Loại: Lý thuyết
% File riêng, không trộn vào de.tex
% Định nghĩa lệnh Lý thuyết — dùng khi biên dịch lt.tex bằng TeX.
% App web đọc lt.tex trực tiếp (bỏ \input file này).
% Trong lt.tex, từ thư mục bài:
%   % Định nghĩa lệnh Lý thuyết — dùng khi biên dịch lt.tex bằng TeX.
% App web đọc lt.tex trực tiếp (bỏ \input file này).
% Trong lt.tex, từ thư mục bài:
%   % Định nghĩa lệnh Lý thuyết — dùng khi biên dịch lt.tex bằng TeX.
% App web đọc lt.tex trực tiếp (bỏ \input file này).
% Trong lt.tex, từ thư mục bài:
%   \input{../../../../_lenh/lythuyet.tex}
% từ gốc repo:
%   \input{ngan-hang/_lenh/lythuyet.tex}
\makeatletter
\providecommand{\indam}[1]{\textbf{#1}}
\providecommand{\chuy}[1]{\par\smallskip\noindent\textit{#1}\par}
\providecommand{\luuy}[1]{\par\smallskip\noindent\textbf{Lưu ý.} #1\par}
\providecommand{\ghichu}[1]{\par\smallskip\noindent\textbf{Ghi chú.} #1\par}
\providecommand{\vidu}[1]{\par\smallskip\noindent\textbf{Ví dụ.} #1\par}
\providecommand{\Blythuyet}{\subsection*{Lý thuyết}}
% Hai cột: kiến thức | hình
\providecommand{\haicotchay}[2]{%
  \par\medskip\noindent
  \begin{minipage}[t]{0.52\linewidth}#1\end{minipage}\hfill
  \begin{minipage}[t]{0.46\linewidth}#2\end{minipage}%
  \par\medskip
}
\providecommand{\kienthuccot}[2]{\haicotchay{#1}{#2}}
% Dự phòng nếu chưa nạp graphicx / caption (không ghi đè bản thật)
\providecommand{\resizebox}[3]{#3}
\providecommand{\captionof}[2]{\par\smallskip\noindent\textit{#2}\par}
\newenvironment{hoatdong}[1]{%
  \par\medskip\noindent\textbf{Hoạt động.} #1\par\smallskip
}{\par\medskip}
\newenvironment{emcobiet}{%
  \par\medskip\noindent\textbf{Em có biết}\par\smallskip
}{\par\medskip}
\newenvironment{emdahoc}{%
  \par\medskip\noindent\textbf{Em đã học}\par\smallskip
}{\par\medskip}
\newenvironment{emcothe}{%
  \par\medskip\noindent\textbf{Em có thể}\par\smallskip
}{\par\medskip}
\newenvironment{kienthuc}{%
  \par\medskip\noindent\textbf{Kiến thức cốt lõi}\par\smallskip
}{\par\medskip}
\newenvironment{traloi}{%
  \par\smallskip\noindent\textit{Trả lời / ghi chú của em}\par
  \vspace{1.2em}%
}{\par\medskip}
\makeatother

% từ gốc repo:
%   % Định nghĩa lệnh Lý thuyết — dùng khi biên dịch lt.tex bằng TeX.
% App web đọc lt.tex trực tiếp (bỏ \input file này).
% Trong lt.tex, từ thư mục bài:
%   \input{../../../../_lenh/lythuyet.tex}
% từ gốc repo:
%   \input{ngan-hang/_lenh/lythuyet.tex}
\makeatletter
\providecommand{\indam}[1]{\textbf{#1}}
\providecommand{\chuy}[1]{\par\smallskip\noindent\textit{#1}\par}
\providecommand{\luuy}[1]{\par\smallskip\noindent\textbf{Lưu ý.} #1\par}
\providecommand{\ghichu}[1]{\par\smallskip\noindent\textbf{Ghi chú.} #1\par}
\providecommand{\vidu}[1]{\par\smallskip\noindent\textbf{Ví dụ.} #1\par}
\providecommand{\Blythuyet}{\subsection*{Lý thuyết}}
% Hai cột: kiến thức | hình
\providecommand{\haicotchay}[2]{%
  \par\medskip\noindent
  \begin{minipage}[t]{0.52\linewidth}#1\end{minipage}\hfill
  \begin{minipage}[t]{0.46\linewidth}#2\end{minipage}%
  \par\medskip
}
\providecommand{\kienthuccot}[2]{\haicotchay{#1}{#2}}
% Dự phòng nếu chưa nạp graphicx / caption (không ghi đè bản thật)
\providecommand{\resizebox}[3]{#3}
\providecommand{\captionof}[2]{\par\smallskip\noindent\textit{#2}\par}
\newenvironment{hoatdong}[1]{%
  \par\medskip\noindent\textbf{Hoạt động.} #1\par\smallskip
}{\par\medskip}
\newenvironment{emcobiet}{%
  \par\medskip\noindent\textbf{Em có biết}\par\smallskip
}{\par\medskip}
\newenvironment{emdahoc}{%
  \par\medskip\noindent\textbf{Em đã học}\par\smallskip
}{\par\medskip}
\newenvironment{emcothe}{%
  \par\medskip\noindent\textbf{Em có thể}\par\smallskip
}{\par\medskip}
\newenvironment{kienthuc}{%
  \par\medskip\noindent\textbf{Kiến thức cốt lõi}\par\smallskip
}{\par\medskip}
\newenvironment{traloi}{%
  \par\smallskip\noindent\textit{Trả lời / ghi chú của em}\par
  \vspace{1.2em}%
}{\par\medskip}
\makeatother

\makeatletter
\providecommand{\indam}[1]{\textbf{#1}}
\providecommand{\chuy}[1]{\par\smallskip\noindent\textit{#1}\par}
\providecommand{\luuy}[1]{\par\smallskip\noindent\textbf{Lưu ý.} #1\par}
\providecommand{\ghichu}[1]{\par\smallskip\noindent\textbf{Ghi chú.} #1\par}
\providecommand{\vidu}[1]{\par\smallskip\noindent\textbf{Ví dụ.} #1\par}
\providecommand{\Blythuyet}{\subsection*{Lý thuyết}}
% Hai cột: kiến thức | hình
\providecommand{\haicotchay}[2]{%
  \par\medskip\noindent
  \begin{minipage}[t]{0.52\linewidth}#1\end{minipage}\hfill
  \begin{minipage}[t]{0.46\linewidth}#2\end{minipage}%
  \par\medskip
}
\providecommand{\kienthuccot}[2]{\haicotchay{#1}{#2}}
% Dự phòng nếu chưa nạp graphicx / caption (không ghi đè bản thật)
\providecommand{\resizebox}[3]{#3}
\providecommand{\captionof}[2]{\par\smallskip\noindent\textit{#2}\par}
\newenvironment{hoatdong}[1]{%
  \par\medskip\noindent\textbf{Hoạt động.} #1\par\smallskip
}{\par\medskip}
\newenvironment{emcobiet}{%
  \par\medskip\noindent\textbf{Em có biết}\par\smallskip
}{\par\medskip}
\newenvironment{emdahoc}{%
  \par\medskip\noindent\textbf{Em đã học}\par\smallskip
}{\par\medskip}
\newenvironment{emcothe}{%
  \par\medskip\noindent\textbf{Em có thể}\par\smallskip
}{\par\medskip}
\newenvironment{kienthuc}{%
  \par\medskip\noindent\textbf{Kiến thức cốt lõi}\par\smallskip
}{\par\medskip}
\newenvironment{traloi}{%
  \par\smallskip\noindent\textit{Trả lời / ghi chú của em}\par
  \vspace{1.2em}%
}{\par\medskip}
\makeatother

% từ gốc repo:
%   % Định nghĩa lệnh Lý thuyết — dùng khi biên dịch lt.tex bằng TeX.
% App web đọc lt.tex trực tiếp (bỏ \input file này).
% Trong lt.tex, từ thư mục bài:
%   % Định nghĩa lệnh Lý thuyết — dùng khi biên dịch lt.tex bằng TeX.
% App web đọc lt.tex trực tiếp (bỏ \input file này).
% Trong lt.tex, từ thư mục bài:
%   \input{../../../../_lenh/lythuyet.tex}
% từ gốc repo:
%   \input{ngan-hang/_lenh/lythuyet.tex}
\makeatletter
\providecommand{\indam}[1]{\textbf{#1}}
\providecommand{\chuy}[1]{\par\smallskip\noindent\textit{#1}\par}
\providecommand{\luuy}[1]{\par\smallskip\noindent\textbf{Lưu ý.} #1\par}
\providecommand{\ghichu}[1]{\par\smallskip\noindent\textbf{Ghi chú.} #1\par}
\providecommand{\vidu}[1]{\par\smallskip\noindent\textbf{Ví dụ.} #1\par}
\providecommand{\Blythuyet}{\subsection*{Lý thuyết}}
% Hai cột: kiến thức | hình
\providecommand{\haicotchay}[2]{%
  \par\medskip\noindent
  \begin{minipage}[t]{0.52\linewidth}#1\end{minipage}\hfill
  \begin{minipage}[t]{0.46\linewidth}#2\end{minipage}%
  \par\medskip
}
\providecommand{\kienthuccot}[2]{\haicotchay{#1}{#2}}
% Dự phòng nếu chưa nạp graphicx / caption (không ghi đè bản thật)
\providecommand{\resizebox}[3]{#3}
\providecommand{\captionof}[2]{\par\smallskip\noindent\textit{#2}\par}
\newenvironment{hoatdong}[1]{%
  \par\medskip\noindent\textbf{Hoạt động.} #1\par\smallskip
}{\par\medskip}
\newenvironment{emcobiet}{%
  \par\medskip\noindent\textbf{Em có biết}\par\smallskip
}{\par\medskip}
\newenvironment{emdahoc}{%
  \par\medskip\noindent\textbf{Em đã học}\par\smallskip
}{\par\medskip}
\newenvironment{emcothe}{%
  \par\medskip\noindent\textbf{Em có thể}\par\smallskip
}{\par\medskip}
\newenvironment{kienthuc}{%
  \par\medskip\noindent\textbf{Kiến thức cốt lõi}\par\smallskip
}{\par\medskip}
\newenvironment{traloi}{%
  \par\smallskip\noindent\textit{Trả lời / ghi chú của em}\par
  \vspace{1.2em}%
}{\par\medskip}
\makeatother

% từ gốc repo:
%   % Định nghĩa lệnh Lý thuyết — dùng khi biên dịch lt.tex bằng TeX.
% App web đọc lt.tex trực tiếp (bỏ \input file này).
% Trong lt.tex, từ thư mục bài:
%   \input{../../../../_lenh/lythuyet.tex}
% từ gốc repo:
%   \input{ngan-hang/_lenh/lythuyet.tex}
\makeatletter
\providecommand{\indam}[1]{\textbf{#1}}
\providecommand{\chuy}[1]{\par\smallskip\noindent\textit{#1}\par}
\providecommand{\luuy}[1]{\par\smallskip\noindent\textbf{Lưu ý.} #1\par}
\providecommand{\ghichu}[1]{\par\smallskip\noindent\textbf{Ghi chú.} #1\par}
\providecommand{\vidu}[1]{\par\smallskip\noindent\textbf{Ví dụ.} #1\par}
\providecommand{\Blythuyet}{\subsection*{Lý thuyết}}
% Hai cột: kiến thức | hình
\providecommand{\haicotchay}[2]{%
  \par\medskip\noindent
  \begin{minipage}[t]{0.52\linewidth}#1\end{minipage}\hfill
  \begin{minipage}[t]{0.46\linewidth}#2\end{minipage}%
  \par\medskip
}
\providecommand{\kienthuccot}[2]{\haicotchay{#1}{#2}}
% Dự phòng nếu chưa nạp graphicx / caption (không ghi đè bản thật)
\providecommand{\resizebox}[3]{#3}
\providecommand{\captionof}[2]{\par\smallskip\noindent\textit{#2}\par}
\newenvironment{hoatdong}[1]{%
  \par\medskip\noindent\textbf{Hoạt động.} #1\par\smallskip
}{\par\medskip}
\newenvironment{emcobiet}{%
  \par\medskip\noindent\textbf{Em có biết}\par\smallskip
}{\par\medskip}
\newenvironment{emdahoc}{%
  \par\medskip\noindent\textbf{Em đã học}\par\smallskip
}{\par\medskip}
\newenvironment{emcothe}{%
  \par\medskip\noindent\textbf{Em có thể}\par\smallskip
}{\par\medskip}
\newenvironment{kienthuc}{%
  \par\medskip\noindent\textbf{Kiến thức cốt lõi}\par\smallskip
}{\par\medskip}
\newenvironment{traloi}{%
  \par\smallskip\noindent\textit{Trả lời / ghi chú của em}\par
  \vspace{1.2em}%
}{\par\medskip}
\makeatother

\makeatletter
\providecommand{\indam}[1]{\textbf{#1}}
\providecommand{\chuy}[1]{\par\smallskip\noindent\textit{#1}\par}
\providecommand{\luuy}[1]{\par\smallskip\noindent\textbf{Lưu ý.} #1\par}
\providecommand{\ghichu}[1]{\par\smallskip\noindent\textbf{Ghi chú.} #1\par}
\providecommand{\vidu}[1]{\par\smallskip\noindent\textbf{Ví dụ.} #1\par}
\providecommand{\Blythuyet}{\subsection*{Lý thuyết}}
% Hai cột: kiến thức | hình
\providecommand{\haicotchay}[2]{%
  \par\medskip\noindent
  \begin{minipage}[t]{0.52\linewidth}#1\end{minipage}\hfill
  \begin{minipage}[t]{0.46\linewidth}#2\end{minipage}%
  \par\medskip
}
\providecommand{\kienthuccot}[2]{\haicotchay{#1}{#2}}
% Dự phòng nếu chưa nạp graphicx / caption (không ghi đè bản thật)
\providecommand{\resizebox}[3]{#3}
\providecommand{\captionof}[2]{\par\smallskip\noindent\textit{#2}\par}
\newenvironment{hoatdong}[1]{%
  \par\medskip\noindent\textbf{Hoạt động.} #1\par\smallskip
}{\par\medskip}
\newenvironment{emcobiet}{%
  \par\medskip\noindent\textbf{Em có biết}\par\smallskip
}{\par\medskip}
\newenvironment{emdahoc}{%
  \par\medskip\noindent\textbf{Em đã học}\par\smallskip
}{\par\medskip}
\newenvironment{emcothe}{%
  \par\medskip\noindent\textbf{Em có thể}\par\smallskip
}{\par\medskip}
\newenvironment{kienthuc}{%
  \par\medskip\noindent\textbf{Kiến thức cốt lõi}\par\smallskip
}{\par\medskip}
\newenvironment{traloi}{%
  \par\smallskip\noindent\textit{Trả lời / ghi chú của em}\par
  \vspace{1.2em}%
}{\par\medskip}
\makeatother

\makeatletter
\providecommand{\indam}[1]{\textbf{#1}}
\providecommand{\chuy}[1]{\par\smallskip\noindent\textit{#1}\par}
\providecommand{\luuy}[1]{\par\smallskip\noindent\textbf{Lưu ý.} #1\par}
\providecommand{\ghichu}[1]{\par\smallskip\noindent\textbf{Ghi chú.} #1\par}
\providecommand{\vidu}[1]{\par\smallskip\noindent\textbf{Ví dụ.} #1\par}
\providecommand{\Blythuyet}{\subsection*{Lý thuyết}}
% Hai cột: kiến thức | hình
\providecommand{\haicotchay}[2]{%
  \par\medskip\noindent
  \begin{minipage}[t]{0.52\linewidth}#1\end{minipage}\hfill
  \begin{minipage}[t]{0.46\linewidth}#2\end{minipage}%
  \par\medskip
}
\providecommand{\kienthuccot}[2]{\haicotchay{#1}{#2}}
% Dự phòng nếu chưa nạp graphicx / caption (không ghi đè bản thật)
\providecommand{\resizebox}[3]{#3}
\providecommand{\captionof}[2]{\par\smallskip\noindent\textit{#2}\par}
\newenvironment{hoatdong}[1]{%
  \par\medskip\noindent\textbf{Hoạt động.} #1\par\smallskip
}{\par\medskip}
\newenvironment{emcobiet}{%
  \par\medskip\noindent\textbf{Em có biết}\par\smallskip
}{\par\medskip}
\newenvironment{emdahoc}{%
  \par\medskip\noindent\textbf{Em đã học}\par\smallskip
}{\par\medskip}
\newenvironment{emcothe}{%
  \par\medskip\noindent\textbf{Em có thể}\par\smallskip
}{\par\medskip}
\newenvironment{kienthuc}{%
  \par\medskip\noindent\textbf{Kiến thức cốt lõi}\par\smallskip
}{\par\medskip}
\newenvironment{traloi}{%
  \par\smallskip\noindent\textit{Trả lời / ghi chú của em}\par
  \vspace{1.2em}%
}{\par\medskip}
\makeatother


Dưới đây là nội dung hoàn chỉnh của file `lt.tex` biên soạn chuẩn theo chương trình SGK Vật lí 12 Kết nối tri thức với cuộc sống (Bài 3: Nhiệt độ. Thang nhiệt độ – nhiệt kế), sử dụng đúng cấu trúc LaTeX, đầy đủ các mục kiến thức và các môi trường yêu cầu.

```latex
%===================================================================================
% BÀI 3: NHIỆT ĐỘ. THANG NHIỆT ĐỘ – NHIỆT KẾ
% Sách giáo khoa Vật lí 12 - Bộ Kết nối tri thức với cuộc sống
%===================================================================================

\subsection{Lý thuyết}

\subsubsection{Khái niệm nhiệt độ và trạng thái cân bằng nhiệt}

\begin{hoatdong}
	\textbf{Tìm hiểu về cảm giác nóng, lạnh và trạng thái cân bằng nhiệt:}
	\begin{enumerate}
		\item Chuẩn bị ba cốc nước: cốc (a) đựng nước lạnh, cốc (b) đựng nước nguội (nhiệt độ phòng), cốc (c) đựng nước ấm.
		\item Nhúng ngón tay trỏ của tay phải vào cốc (a), ngón tay trỏ của tay trái vào cốc (c). Sau một lúc, rút cả hai ngón tay ra rồi cùng nhúng vào cốc (b). Nêu cảm giác của em về độ nóng, lạnh ở mỗi ngón tay. Cảm giác của tay có thể xác định chính xác mức độ nóng, lạnh của một vật hay không?
		\item Khi thả một thỏi kim loại đang nóng vào một chậu nước lạnh, sau một thời gian nhiệt độ của thỏi kim loại và nước thay đổi như thế nào?
	\end{enumerate}
\end{hoatdong}

\begin{kienthuc}
	\begin{enumerate}
		\item \textbf{Khái niệm nhiệt độ:}
		\begin{itemize}
			\item \textit{Ở góc độ vĩ mô:} Nhiệt độ là đại lượng vật lý đặc trưng cho mức độ chuyển động nhiệt hỗn loạn của các phân tử cấu tạo nên vật và cho biết mức độ nóng hay lạnh của vật đó. Nhiệt độ quyết định chiều tự truyền nhiệt lượng giữa hai vật khi cho tiếp xúc nhiệt với nhau: nhiệt lượng luôn tự truyền từ vật có nhiệt độ cao sang vật có nhiệt độ thấp hơn.
			\item \textit{Ở góc độ vi mô:} Nhiệt độ là số đo động năng tịnh tiến trung bình của các phân tử cấu tạo nên vật. Khi vật có nhiệt độ càng cao thì các phân tử chuyển động hỗn loạn càng nhanh, dẫn đến động năng tịnh tiến trung bình của chúng càng lớn:
			\[
			\bar{E}_{\mathrm{d}} = \dfrac{3}{2} k T
			\]
			trong đó $k \approx 1{,}38 \cdot 10^{-23}\ \mathrm{J/K}$ là hằng số Boltzmann, $T$ là nhiệt độ tuyệt đối của vật.
		\end{itemize}
		\item \textbf{Trạng thái cân bằng nhiệt:}
		\begin{itemize}
			\item Khi hai vật có nhiệt độ khác nhau tiếp xúc nhiệt với nhau, nhiệt năng sẽ truyền từ vật có nhiệt độ cao hơn sang vật có nhiệt độ thấp hơn cho đến khi nhiệt độ của hai vật bằng nhau. Trạng thái này được gọi là \textbf{trạng thái cân bằng nhiệt}.
			\item Khi đạt trạng thái cân bằng nhiệt, không còn sự truyền nhiệt ròng giữa hai vật, và động năng tịnh tiến trung bình của các phân tử trong hai vật là như nhau.
		\end{itemize}
	\end{enumerate}
\end{kienthuc}

\chuy{Nhiệt độ là đại lượng đặc trưng cho động năng \textbf{trung bình} của mỗi phân tử trong vật (đại lượng cường tính), hoàn toàn không phải là tổng nhiệt năng (tổng động năng chuyển động nhiệt của toàn bộ các phân tử) của vật. Một giọt nước sôi ở $100^\circ\mathrm{C}$ có nhiệt độ cao hơn một bể bơi chứa hàng ngàn mét khối nước ở $25^\circ\mathrm{C}$, nhưng tổng nhiệt năng của bể bơi lớn hơn rất nhiều so với giọt nước.}

\subsubsection{Thang nhiệt độ}

\begin{kienthuc}
	Để đo nhiệt độ, người ta phải định ra thang đo với các điểm mốc chuẩn cố định:
	\begin{enumerate}
		\item \textbf{Thang nhiệt độ Celsius ($^\circ\mathrm{C}$):}
		\begin{itemize}
			\item Chọn mốc $0^\circ\mathrm{C}$ là nhiệt độ đóng băng của nước tinh khiết (nước đá đang tan) ở áp suất tiêu chuẩn ($1\ \mathrm{atm}$).
			\item Chọn mốc $100^\circ\mathrm{C}$ là nhiệt độ sôi của nước tinh khiết (hơi nước đang sôi) ở áp suất tiêu chuẩn.
			\item Khoảng giữa hai mốc này được chia đều thành $100$ phần bằng nhau, mỗi phần tương ứng với $1^\circ\mathrm{C}$.
		\end{itemize}
		\item \textbf{Thang nhiệt độ Kelvin (Thang nhiệt độ tuyệt đối - $\mathrm{K}$):}
		\begin{itemize}
			\item \textbf{Độ không tuyệt đối ($0\ \mathrm{K}$):} Là mức nhiệt độ thấp nhất trên lý thuyết mà vật chất có thể đạt tới. Tại $0\ \mathrm{K}$, toàn bộ chuyển động nhiệt hỗn loạn của các phân tử bị triệt tiêu hoàn toàn (động năng phân tử bằng không). Mọi vật chất trong vũ trụ \textbf{không thể} có nhiệt độ thấp hơn $0\ \mathrm{K}$.
			\item Giá trị $0\ \mathrm{K}$ tương ứng với $-273{,}15^\circ\mathrm{C}$ (trong tính toán thông thường lấy gần đúng là $-273^\circ\mathrm{C}$).
			\item Mỗi độ chia trong thang Kelvin có độ lớn đúng bằng một độ chia trong thang Celsius: $1\ \mathrm{K} = 1^\circ\mathrm{C}$.
			\item Hệ thức chuyển đổi giữa hai thang đo:
			\[
			T\ (\mathrm{K}) = t\ (^\circ\mathrm{C}) + 273{,}15 \quad (\text{hay } T \approx t + 273)
			\]
			\item Độ biến thiên nhiệt độ trong hai thang đo là bằng nhau:
			\[
			\Delta T\ (\mathrm{K}) = \Delta t\ (^\circ\mathrm{C})
			\]
		\end{itemize}
		\item \textbf{Thang nhiệt độ Fahrenheit ($^\circ\mathrm{F}$):}
		\begin{itemize}
			\item Được sử dụng phổ biến ở các nước nói tiếng Anh (Mỹ, Anh...). Mốc nước đá đang tan là $32^\circ\mathrm{F}$, mốc nước đang sôi là $212^\circ\mathrm{F}$.
			\item Khoảng giữa hai mốc được chia thành $180$ phần bằng nhau, mỗi phần tương ứng với $1^\circ\mathrm{F}$.
			\item Công thức chuyển đổi qua lại giữa thang Celsius và thang Fahrenheit:
			\[
			t\ (^\circ\mathrm{F}) = 1{,}8 \cdot t\ (^\circ\mathrm{C}) + 32 \quad \Longleftrightarrow \quad t\ (^\circ\mathrm{C}) = \dfrac{t\ (^\circ\mathrm{F}) - 32}{1{,}8}
			\]
		\end{itemize}
	\end{enumerate}
\end{kienthuc}

\begin{emcobiet}
	\textbf{Quy tắc chuyển đổi tổng quát cho thang nhiệt độ tự chế:}
	
	Nếu gọi một thang nhiệt độ tự chế là thang $\mathrm{Z}$, trong đó mốc nhiệt độ nước đá đang tan là $t_{0\mathrm{Z}}$ và mốc nước đang sôi ở áp suất tiêu chuẩn là $t_{100\mathrm{Z}}$, do sự biến thiên giữa các thang đo là tuyến tính bậc nhất, ta luôn có tỉ số tương đương:
	\[
	\dfrac{t_\mathrm{C} - 0}{100 - 0} = \dfrac{t_\mathrm{Z} - t_{0\mathrm{Z}}}{t_{100\mathrm{Z}} - t_{0\mathrm{Z}}} = \dfrac{T - 273{,}15}{100} = \dfrac{t_\mathrm{F} - 32}{212 - 32}
	\]
	Quy tắc tỷ lệ này cho phép thiết lập và chuyển đổi số chỉ giữa mọi thang đo nhiệt độ với nhau một cách chính xác.
\end{emcobiet}

\subsubsection{Nhiệt kế và nguyên lý hoạt động}

\begin{kienthuc}
	\begin{enumerate}
		\item \textbf{Nguyên lý hoạt động của nhiệt kế:}
		Nhiệt kế là dụng cụ dùng để đo nhiệt độ. Hoạt động của nhiệt kế dựa trên việc đo một đại lượng vật lý phụ thuộc tuyến tính hoặc liên tục vào nhiệt độ (gọi là đại lượng đo nhiệt độ).
		\begin{itemize}
			\item \textbf{Nhiệt kế chất lỏng (thủy ngân, rượu, dầu...):} Hoạt động dựa trên \textbf{sự nở vì nhiệt của chất lỏng}. Khi nhiệt độ tăng, chất lỏng trong bầu dâng lên trong ống quản hẹp có thang chia độ.
			\item \textbf{Nhiệt kế kim loại (băng kép):} Hoạt động dựa trên sự nở vì nhiệt khác nhau của hai thanh kim loại ghép dính vào nhau làm biến dạng cong khi nhiệt độ thay đổi.
			\item \textbf{Nhiệt kế điện trở:} Hoạt động dựa trên sự thay đổi điện trở của dây kim loại (như bạch kim - Pt) hoặc chất bán dẫn (nhiệt điện trở) theo nhiệt độ.
			\item \textbf{Nhiệt kế hồng ngoại:} Hoạt động dựa trên việc đo cường độ bức xạ hồng ngoại do vật nóng phát ra, giúp đo nhiệt độ mà không cần tiếp xúc trực tiếp (thường dùng đo thân nhiệt nhanh hoặc đo vật thể nóng từ xa).
		\end{itemize}
		\item \textbf{Mối liên hệ giữa chiều cao cột chất lỏng và nhiệt độ:}
		Ở nhiệt kế chất lỏng có ống quản tiết diện đều, chiều cao cột chất lỏng $h$ phụ thuộc tuyến tính vào nhiệt độ $t$:
		\[
		\dfrac{t - t_1}{t_2 - t_1} = \dfrac{h - h_1}{h_2 - h_1}
		\]
		trong đó $h_1, h_2$ lần lượt là chiều cao cột chất lỏng tương ứng với các nhiệt độ mốc $t_1, t_2$.
	\end{enumerate}
\end{kienthuc}

\chuy{
	\begin{itemize}
		\item Không thể dùng nhiệt kế thủy ngân hay nhiệt kế rượu để đo nhiệt độ cực thấp gần độ không tuyệt đối ($0\ \mathrm{K}$) vì thủy ngân đông đặc ở $-38{,}83^\circ\mathrm{C}$ và rượu ethanol đông đặc ở $-114{,}1^\circ\mathrm{C}$.
		\item Khi sử dụng nhiệt kế thủy ngân bị vỡ cần hết sức cẩn thận, không dùng tay trần thu gom trực tiếp vì hơi thủy ngân rất độc; có thể dùng bột lưu huỳnh rắc lên để chuyển hóa thủy ngân lỏng thành hợp chất $\mathrm{HgS}$ an toàn hơn trước khi dọn dẹp.
	\end{itemize}
}

\subsubsection{Các cách làm biến đổi nội năng của một vật}

\begin{kienthuc}
	Nội năng của một vật là tổng động năng của chuyển động nhiệt của các phân tử và thế năng tương tác giữa các phân tử cấu tạo nên vật: $U = E_{\mathrm{d\Sigma}} + E_{\mathrm{t\Sigma}}$. Có hai cách làm biến đổi nội năng:
	\begin{enumerate}
		\item \textbf{Thực hiện công (Công cơ học):}
		\begin{itemize}
			\item Là quá trình làm thay đổi nội năng có sự tham gia của các lực cơ học tác dụng và sinh công làm chuyển dời điểm đặt của lực.
			\item Trong quá trình này có sự chuyển hóa từ một dạng năng lượng khác (như cơ năng) thành nội năng của vật (hoặc ngược lại).
			\item \textit{Ví dụ:} Cọ xát hai bàn tay vào nhau thấy tay nóng lên; nén khí trong ống bơm xe đạp làm pit-tông và không khí bên trong nóng lên.
		\end{itemize}
		\item \textbf{Truyền nhiệt (Nhiệt lượng):}
		\begin{itemize}
			\item Là quá trình làm thay đổi nội năng của vật mà \textbf{không có sự thực hiện công} cơ học (không có sự chuyển dời vĩ mô).
			\item Trong quá trình này không có sự chuyển hóa năng lượng giữa các dạng khác nhau, mà chỉ có sự truyền trực tiếp nội năng từ vật này sang vật khác (hoặc từ phần này sang phần khác của cùng một vật) do có sự chênh lệch nhiệt độ.
			\item Số đo độ biến thiên nội năng trong quá trình truyền nhiệt được gọi là \textbf{nhiệt lượng} ($Q$):
			\[
			\Delta U = Q
			\]
			\item \textit{Ví dụ:} Thả một thìa kim loại nguội vào cốc nước sôi thì thìa nóng dần lên; phơi thanh sắt ngoài trời nắng.
		\end{itemize}
	\end{enumerate}
\end{kienthuc}

\begin{emdahoc}
	\begin{itemize}
		\item Nhiệt độ phản ánh động năng tịnh tiến trung bình của các phân tử cấu tạo nên vật: $\bar{E}_{\mathrm{d}} = \dfrac{3}{2} k T$.
		\item Độ không tuyệt đối ($0\ \mathrm{K} = -273{,}15^\circ\mathrm{C}$) là mức nhiệt độ thấp nhất trên lý thuyết, không vật nào có thể có nhiệt độ dưới $0\ \mathrm{K}$.
		\item Chuyển đổi giữa ba thang nhiệt độ thông dụng:
		\[
		T\ (\mathrm{K}) = t\ (^\circ\mathrm{C}) + 273{,}15 \quad ; \quad t\ (^\circ\mathrm{F}) = 1{,}8 \cdot t\ (^\circ\mathrm{C}) + 32
		\]
		\item Nhiệt kế hoạt động dựa vào việc đo sự thay đổi của một đại lượng vật lý theo nhiệt độ (như sự nở vì nhiệt của chất lỏng, điện trở kim loại...).
		\item Có hai cách làm biến đổi nội năng của vật là \textbf{thực hiện công} (có sự chuyển hóa cơ năng thành nội năng) và \textbf{truyền nhiệt} (chỉ truyền nội năng, không có công cơ học).
	\end{itemize}
\end{emdahoc}

\begin{emcothe}
	\begin{itemize}
		\item Giải thích được nguyên lý đo thân nhiệt bằng súng nhiệt kế hồng ngoại dùng trong y tế và phòng chống dịch bệnh.
		\item Tự thiết kế và chia vạch một nhiệt kế chất lỏng tự chế đơn giản bằng các mốc chuẩn nhiệt độ quen thuộc (nước đá tan $0^\circ\mathrm{C}$, nước sôi $100^\circ\mathrm{C}$).
		\item Hiệu chỉnh lại số đo khi dùng một nhiệt kế bị vạch chia sai lệch mốc chuẩn bằng phương pháp nội suy tuyến tính.
	\end{itemize}
\end{emcothe}
